\homework{9}{Tue 04 Apr 2023 19:18}{}

45.B, 45.C, 45.D, 45.G, 45.K, 45.M, 45.N

\begin{itemize}
	\item 45.B. Consider the polar coordinate $\operatorname{map}(x, y)=\varphi(r, \theta)=(r \cos \theta, r \sin \theta)$ defined on $\boldsymbol{R}^2$, and its behavior on the set $A=[0,1] \times[0,2 \pi]$. Use Theorem 45.4 to obtain the reassuring information that the image $D=\varphi(A)$, which is the unit disk $D=\left\{(x, y): x^2+y^2 \leq 1\right\}$, has content. Investigate the manner in which $\varphi$ maps the boundary of $A$. Show that the boundary of $D$ is the image under $\varphi$ of only one side of $A$, and that the other three sides of $A$ get mapped into the interior of $D$.
\begin{proof}
	Clearly $A$ has content, and for all $x \in \text{int} A$, we have $x \in (0,1) \times (0, 2 \pi)$. Thus
	\[
		J_\varphi(x) = 
		\begin{vmatrix}
			\cos \theta & r \sin \theta \\
			\sin \theta & - r \cos \theta \\
		\end{vmatrix}
		= - r \cos^2 \theta - r \sin^2 \theta = - r \neq 0
	.\] 
	Boundary of $D$ is
	 \[
		\partial D = \{(x,y): x^2 + y^2 = 1\}
	.\] 
	Where
	\[
		\varphi\left( \{1\} \times [0,2 \pi] \right) = \{(\cos \theta, \sin \theta): \theta \in [0,2 \pi]\} = \partial D
	.\] 
	For the other three boundaries, $r < 1$. Thus $x^2 + y^2 = r^2 < 1$. Hence $\varphi(r, \theta)$ is in the interior of $D$.
\end{proof}


	\item 45.C. Consider the map $(x, y)=\psi(u, v)=(\sin u, \sin v)$ defined on $\boldsymbol{R}^2$. Determine the image of the boundary of $B=\left[-{ }_4^3 \pi,{ }_4^3 \pi\right] \times\left[-\frac{3}{4} \pi, \frac{3}{4} \pi\right]$ under $\psi$, and the boundary of $\psi(B)$. Show that most, though not quite all, of the boundary points of $\psi(B)$ are images of interior points of $B$.
\begin{proof}
\end{proof}


	\item 45.D. Given that the area of the circular disk $\left\{(x, y): x^2+y^2 \leq 1\right\}$ is equal to $\pi$, find the areas of the elliptical disks given by:
(a) $\left\{(x, y): \frac{x^2}{4}+\frac{y^2}{9} \leq 1\right\}$
(b) $\left\{(x, y): 2 x^2+2 x y+5 y^2 \leq 1\right\}$.
(Hint: $\left.2 x^2+2 x y+5 y^2=(x+2 y)^2+(x-y)^2.\right)$
\item 45.G. Evaluate the iterated integral
\[
\int_1^3\left\{\int_{x^2}^{x^2+1} x y d y\right\} d x
\]
directly. Then use the transformation $(x, y) \mapsto(u, v)=\left(x, y-x^2\right)$ to evaluate this integral.

\item 45.K. Let $\alpha<\beta$ belong to $[0,2 \pi]$ and let $h:[\alpha, \beta] \rightarrow \mathbf{R}$ be continuous and such that $h(\theta) \geq 0$ for $\theta \in[\alpha, \beta]$. Let $H=\left\{(\theta, r) \in \boldsymbol{R}^2: \alpha \leq \theta \leq \beta, 0 \leq r \leq h(\theta)\right\}$ be the ordinate set of $h$ (see Exercise $44 . O$ ), so that $H$ has content. The polar curve generated by $h$ is the curve in $\boldsymbol{R}^2$ defined by $\theta \mapsto(h(\theta) \cos \theta, h(\theta) \sin \theta)$, and the polar ordinate set of this curve is the set
\[
H_1=\left\{(r \cos \theta, r \sin \theta) \in \mathbf{R}^2: \alpha \leq \theta \leq \beta, 0 \leq r \leq h(\theta)\right\} .
\]
Note that $H_1$ is the image of $H$ under the (reversed) polar map $\varphi_1(\theta, r)=$ $(r \cos \theta, r \sin \theta)$ and use Theorem 45.11 to show that
\[
c\left(H_1\right)=\frac{1}{2} \int_\alpha^\beta(f(\theta))^2 d \theta .
\]
\item 45.M. Let $0 \leq a<b$ and let $f:[a, b] \rightarrow \boldsymbol{R}$ and $S_f$ be as in the preceding exercise. Let $\rho_y: \mathbf{R}^3 \rightarrow \boldsymbol{R}^3$ be defined by $\rho_y(x, y, \theta)=(x \cos \theta, y, x \sin \theta)$ and let $Y_f$ be the image of $S_f \times[0,2 \pi]$ under $\rho_y$. (The set $Y_f$ is called the "solid of revolution obtained by revolving the ordinate set $S_f$ about the $y$-axis.") Use Theorem 45.11 to show that
\[
c\left(Y_f\right)=2 \pi \int_a^b x f(x) d x
\]
\item 45.N. (a) By changing to polar coordinates, show that
\[
\iint_{C_R} e^{-\left(x^2+y^2\right)} d(x, y)=\frac{\pi}{4}\left(1-e^{-R^2}\right),
\]
where $C_R=\left\{(x, y): 0 \leq x, 0 \leq y, x^2+y^2 \leq R^2\right\}$.
(b) If $B_L=\{(x, y): 0 \leq x \leq L, 0 \leq y \leq L\}$, show that
\[
\iint_{B_L} e^{-\left(x^2+y^2\right)} d(x, y)=\left(\int_0^L e^{-x^2} d x\right)^2 .
\]
(c) From the fact that $C_R \subseteq B_R \subseteq C_R \sqrt{2}$, show that
\[
\lim _{R \rightarrow \infty}\left(\int_0^R e^{-x^2} d x\right)^2=\frac{\pi}{4}
\]
whence it follows that $\int_0^{+\infty} e^{-x^2} d x=\frac{1}{2} \sqrt{\pi}$
\end{itemize}
